% Page 013 — page imprimée 9
% Jean Roman, le prédicant.
% Reconstruction éditable en LaTeX à partir du scan.



{\large\bfseries Jean Roman, le prédicant.}

\vspace{0.35em}

Né en 1667 à Vercheny, dans la vallée de la Drôme, dans « une bonne et honnête famille » que le pasteur Julien, réfugié à Lausanne disait avoir connue à la Motte Chalancon, village de la montagne drômoise. Avec sa famille il avait quitté le royaume en 1685, à la Révocation de l’Edit de Nantes, mais revint en France pour prêcher l’Evangile et encourager ses frères en la foi, non pas dans la Drôme où il craignait d’être dissuadé par ses proches, mais dans les Cévennes. Il n’avait pas 20 ans !! Il se déplaçait en tant que colporteur d’où son surnom de « Marchandou ».

Il présida à Meyrueis et dans les environs de nombreuses assemblées dont nous avons la trace par un interrogatoire que le subdélégué Daudé fit à un certain Poujol de Pourcarès ; En voici la teneur :

\vspace{0.35em}

\begin{tcolorbox}[
  colback=white,
  colframe=black,
  boxrule=0.7pt,
  arc=0pt,
  left=0.20cm,
  right=0.20cm,
  top=0.10cm,
  bottom=0.10cm,
  boxsep=0pt
]
\small

« Au nord de l’Aigoual, le subdélégué Daudé s’efforçait de surveiller Roman. Le Marchandou avait, aux environs de Meyrueis, de bons amis. Le jour de la Saint Jean (24 Juin) Roman devant prêcher au valat de Graviès dans le bois des Oubrets, Pierre Bertrand, de Campis, son conducteur s’empressa de convoquer un certain Poujol du hameau voisin de Pourcarès. Poujol déclina violemment l’invitation et traita l’avertisseur de misérable qui ferait périr tout le lieu. Bertrand se borna à lui dire « qu’il s’en laissat ! Il lui faisait trop d’honneur de l’avertir! Désormais il ne lui parlerait de sa vie ! »

Le soir même les soldats de Meyrueis arrêtèrent quelques personnes qui revenaient du culte. Une information fut ouverte qui recueillit de la bouche de Poujol d’utiles renseignements. La maison de Bertrand était le refuge ordinaire de Roman, de frère David et de Plan. Les prédicants y faisaient blanchir leur linge. Un cache y protégeait des armes et des sermons. Bertrand de l’aveu de son père vivait maritalement avec Marie Carteirade de Ferrussac qu’il avait épousée sans autre cérémonie que la bénédiction nuptiale prononcée par Roman en présence de sept témoins.

Tandis que Daudé notait ces détails Bertrand se présenta chez Poujol qui venait de se marier mais en passant, lui, devant le prêtre et vint dire à sa femme « qu’elle était une misérable, de s’être mariée dans une méchante religion et que son mari était un vendeur de chrétiens ».

A son retour du Vigan, Poujol, à son tour fut vertement accueilli par le maire le Seigneur de Caunelles : « Eh bien ! On vous a argenté ? vous êtes un bougre. Vous ne valez pas grand argent, mais contentez vous de le manger et de vous taire ! ». Entre temps Pierre Plan, venu vers Meyrueis, avait encore prêché (début d’août) devant 80 auditeurs dans un ruisseau situé entre Campis et le Crouzet.

Roman se dirige vers Florac en Juillet Août 1696 puis revint encore à Meyrueis. Mais la traîtrise de Poujol dans ce quartier qui lui était familier, constituait pour lui une perpétuelle menace. Bertrand pour réduire le faux frère au silence, usa tout à la fois de l’intimidation et de la prière. Vers la fin Août en pleine nuit, il alla frapper à sa porte suivi du prédicant et de deux amis. Poujol se leva, fit entrer les quatre hommes qui revenaient du Villaret (d’une assemblée sans doute), et dut entendre une longue exhortation de Roman qui le supplia de demeurer ferme, et lui déclara « qu’il ne le croyait pas homme à les vendre ».

Un conflit d’intérêt séparait également Poujol de Bertrand. Bertrand voulut que Roman décide entre eux, mais celui-ci se récusa, « n’étant pas homme d’affaires ». Poujol allait se taire toute une année ; mais dans la crainte d’une surprise Roman tint désormais ses assemblées sous la garde d’hommes armés. Une assemblée pieuse tenue dans la grange du pré de Ribevenès fut protégée par cinq ou six fusiliers postés à la porte… ».

\begin{flushright}
Charles Bost : Les Prédicants Protestants\\[-0.1em]
{\scriptsize p.159}
\end{flushright}

\end{tcolorbox}

